\begin{table}[t]
\centering
\caption{Secondary paired moving-block-bootstrap uncertainty intervals for the frozen PJM test improvements. Circular blocks of 14 days and 20,000 resamples are used; 7- and 28-day sensitivity results are provided in the artifact. The analysis was added after the frozen test outcomes and does not alter the decisions.}
\label{tab:pjm_bootstrap_ci}
\footnotesize
\setlength{\tabcolsep}{3pt}
\begin{tabular}{lll}
\toprule
Region & RMSE improvement [95\% CI] & Cost improvement [95\% CI] \\
\midrule
AEP & 57.97\% [53.06, 61.61] & 58.64\% [54.00, 62.43] \\
COMED & 55.40\% [49.58, 59.83] & 54.11\% [48.31, 58.81] \\
DAYTON & 49.13\% [43.81, 53.53] & 50.30\% [45.16, 54.74] \\
PJME & 51.70\% [45.80, 56.29] & 51.42\% [45.71, 56.34] \\
\bottomrule
\end{tabular}
\end{table}
